\documentclass[11pt,letterpaper]{article}
\usepackage{fontspec}
\usepackage{tocloft}
\usepackage{multicol}
\usepackage{fancyhdr}
\usepackage{nicefrac}
\usepackage[left=1.3in,right=1.3in,top=1in,bottom=1in]{geometry}
\setmainfont[Scale=1.3,AutoFakeBold=1.5,AutoFakeSlant=0.3]{Doves Type}

\widowpenalty=10000
\clubpenalty=10000
\interlinepenalty=500

\title{Chicken Tinga}
\author{}
\date{}

% Configure fancy header
\pagestyle{fancy}
\fancyhf{} % Clear all header and footer fields
\fancyhead[C]{\textit{Chicken Tinga}} % Center header with recipe title
\fancyfoot[C]{\thepage} % Center page number at bottom
\renewcommand{\headrulewidth}{0pt} % Remove header line
\renewcommand{\footrulewidth}{0pt} % Remove footer line

\begin{document}

\maketitle
\thispagestyle{empty}

% \begin{quote}\textit{Bone-in chicken leg quarters are pressure-cooked with fire-roasted tomatoes, a whole can of chipotles in adobo, onion, garlic, and Mexican spice blend. The sauce is blended smooth in the pot, reduced until it coats the meat, and finished with lime and cilantro. One 6-quart Instant Pot batch.}\end{quote}

\section*{Ingredients}
\setlength{\columnsep}{20pt}
\begin{multicols}{2}
\noindent
    Lard \dotfill 2~Tbsp. \\
    White onion, diced \dotfill 1 medium \\
    Garlic cloves, minced \dotfill 4 \\
    Mexican spice blend \dotfill 2~Tbsp. \\
    Chicken leg quarters \dotfill 4~lb. \\
    Diced tomatoes \dotfill 1 (14.5~oz.) can \\
    \columnbreak
    Chipotle chiles in adobo \dotfill 1 (8~oz.) can \\
    Chicken broth \dotfill 1~cup \\
    Bay leaf \dotfill 1 \\
    Lime, juiced \dotfill 1 \\
    Kosher salt \dotfill to taste \\
    Fresh cilantro, chopped \dotfill \nicefrac{1}{4}~cup \\
\end{multicols}

\section*{Directions}

\noindent
Dice the \textbf{white~onion} ---
Mince the \textbf{garlic} ---
Measure the \textbf{Mexican~spice~blend} ---
Open the \textbf{tomatoes} and the \textbf{chipotle~chiles} ---
Measure the \textbf{chicken~broth} ---
Juice the \textbf{lime} ---
Chop the \textbf{cilantro}

\begin{enumerate}
    \item Set a 6-quart Instant Pot to Sauté (Normal). Add 2~Tbsp. \textbf{lard} and heat until it shimmers. Add the \textbf{white~onion} and cook, stirring, for \textit{5--7~minutes}, until the pieces are translucent and soft, the edges just take on color, and the pot looks glossy rather than watery. Continue cooking for \textit{1--2~minutes} if the onion is still crisp or sitting in liquid.

    \item Add the \textbf{garlic} and 2~Tbsp. \textbf{Mexican~spice~blend}. Stir constantly for \textit{30--45~seconds}, until the mixture smells fragrant and the raw spice smell turns toasty.

    \item Nestle in the \textbf{leg~quarters}. Add the \textbf{tomatoes}, the \textbf{chipotle~chiles} with their sauce, 1~cup \textbf{chicken~broth}, and the \textbf{bay~leaf}. Scrape the bottom so the fond dissolves into the liquid. The liquid should cover the pot floor and come partway up the \textbf{chicken}. Cancel Sauté.

    \item Seal the lid and set the valve to Sealing. Cook on high pressure for \textit{15~minutes}. Let the pressure fall on its own for \textit{10~minutes}, then turn the valve to Venting and release any pressure that remains. The \textbf{chicken} is done when the thickest part near the bone reads at least \textit{165°F}.

    \item Transfer the \textbf{chicken} to a cutting board with tongs.

    \item Discard the \textbf{bay~leaf}. Blend the sauce in the pot with an immersion blender until it is smooth and brick-red, with no pieces of chile, tomato, or onion left, about \textit{30--60~seconds}. Keep the blender head under the surface so the hot sauce stays in the pot. Blend for another \textit{15--20~seconds} if any pieces remain.

    \item Set the pot to Sauté (Normal) and simmer the sauce uncovered, stirring often, for \textit{8--15~minutes}, until it coats the back of a spoon. The surface looks glossy, and the sauce still stirs easily. Stir so the bottom does not catch. While it simmers, pull the skin, bones, and any gristle off the \textbf{chicken} and shred the meat with two forks into coarse strands.

    \item Ladle just enough sauce onto the shredded chicken to thoroughly coat it, mixing until every strand is glossy and saucy but not swimming. Reserve and refrigerate or freeze any extra sauce for a future batch or use as a smoky, spicy condiment.

    \item Cancel Sauté. Stir in the \textbf{lime~juice}. Taste. The adobo has already salted the sauce. Add \textbf{kosher~salt} only if the sauce tastes flat. Scatter the \textbf{cilantro} over the top. The tinga is at its best as soon as the \textbf{cilantro} goes on. It holds warm in the pot for about \textit{30~minutes}. Stir in a spoonful of water if it tightens while it waits.
\end{enumerate}

\newpage

% Begin compact two-column layout
{\small
\setlength{\columnsep}{20pt}
\setlength{\multicolsep}{6pt}
\begin{multicols}{2}
\setlength{\parindent}{0pt}
\setlength{\parskip}{4pt}

\section*{Equipment Required}
\begin{itemize}
    \item 6-quart or larger Instant Pot
    \item Immersion blender
    \item Tongs
    \item Two forks
    \item Cutting board
    \item Chef's knife
    \item Wooden spoon or spatula
    \item Measuring cups and spoons
    \item Can opener
    \item Instant-read thermometer
\end{itemize}

\section*{Hints and Notes}

\subsection*{Yield}
\begin{itemize}
    \item One 6-quart Instant Pot batch from 3\nicefrac{1}{2}--4~lb. bone-in \textbf{leg~quarters}.
\end{itemize}

\subsection*{Mise en Place}
\begin{itemize}
    \item Pull the skin off the \textbf{leg~quarters} and trim loose fat
    \item Dice the \textbf{white~onion} and mince the \textbf{garlic}
    \item Measure 2~Tbsp. \textbf{Mexican~spice~blend} and 1~cup \textbf{chicken~broth}
    \item Open the \textbf{tomatoes} and the can of \textbf{chipotle~chiles}
    \item Juice the \textbf{lime} and chop the \textbf{cilantro} so both are ready at the end
    \item Set the immersion blender, tongs, and cutting board beside the pot before you seal the lid
\end{itemize}

\subsection*{Ingredient Tips}
\begin{itemize}
    \item \textbf{Leg~quarters}: bone-in, skin off. Go by weight. About four quarters land in the 3\nicefrac{1}{2}--4~lb. range. The bones and the long cook give the sauce body.
    \item \textbf{Chipotle~chiles}: use the entire 8~oz. can, peppers and sauce. Many grocery cans are labeled 7~oz. Use that whole can the same way. A full can is hot, smoky, and salty.
    \item \textbf{Mexican~spice~blend}: use the blend in this book. It brings cumin, oregano, and a little clove and cinnamon under the chipotle.
    \item \textbf{Lard}: traditional, and it softens the chile heat. Vegetable oil or corn oil works in its place.
    \item \textbf{Tomatoes}: fire-roasted diced tomatoes add char. Plain diced tomatoes work, with less roast flavor.
    \item \textbf{Chicken~broth}: use low-sodium broth. The adobo already seasons the pot. The cup of broth lets the pot reach pressure. The later simmer is what makes the sauce thick.
    \item \textbf{White~onion}: the usual onion for tinga. A yellow onion is fine.
\end{itemize}

\subsection*{Preparation Tips}
\begin{itemize}
    \item \textbf{Pot~fill}: stack the \textbf{leg~quarters} so the food stays under the two-thirds line. An 8-quart pot gives more room for 4~lb. of bone-in chicken.
    \item \textbf{Onion}: let the onion liquid cook off before the \textbf{garlic} and spices go in. Spices bloom in fat. They taste raw if they hit a wet pot.
    \item \textbf{Bay~leaf}: take it out before blending. A blended leaf leaves bitter flecks in the sauce.
    \item \textbf{Blender}: keep the head under the surface and pulse. A head above the sauce throws hot liquid.
    \item \textbf{Reduction}: the \textbf{chicken} releases a lot of juice under pressure, and the can adds more liquid. Judge the sauce by the spoon test, not by the clock.
    \item \textbf{Salt}: taste after the \textbf{lime}. Add \textbf{kosher~salt} only if the sauce tastes flat. Reducing the sauce makes the adobo salt stronger.
    \item \textbf{Shred}: coarse strands hold up in a taco. Paste-fine shreds turn soft in the sauce.
    \item \textbf{Fat}: leave the rendered fat in the sauce. It carries the chipotle and rounds the heat.
\end{itemize}

\subsection*{Make Ahead \& Storage}
\begin{itemize}
    \item The tinga is better the next day. Cool it, then refrigerate up to \textit{4~days}.
    \item Freeze up to \textit{3~months}. Thaw overnight in the refrigerator.
    \item Reheat on the stovetop over low heat. Add a splash of \textbf{broth} if the sauce has tightened.
    \item Add the \textbf{cilantro} after reheating so it stays green.
    \item The best moment is right after the \textbf{lime} and \textbf{cilantro} go on. Warm tinga holds in the pot for about \textit{30~minutes}.
\end{itemize}

\subsection*{Serving Suggestions}
\begin{itemize}
    \item Spoon into warm corn tortillas, onto tostadas, or into tortas.
    \item Top with diced \textbf{white~onion}, avocado, queso fresco, and a squeeze of \textbf{lime}.
    \item Serve with Mexican rice and beans.
    \item A Mexican lager or a hibiscus agua fresca stands up to the chile.
\end{itemize}

\end{multicols}
}

\end{document}
